% 钟面 3:20 时分针夹角 20°
% 12 在上方 (90°)；每个大格 30°，顺时针。
% 3:20：分针 → 4 的位置（数字 4 在 12 顺时针 120° → 极角 90-120=-30°）
%       时针在 3 处 + 10°（已走过 100°，极角 90-100=-10°）
\documentclass[tikz,border=4pt]{standalone}
\usepackage{ctex}
\usepackage{amsmath}
\usetikzlibrary{calc, arrows.meta}
\begin{document}
\begin{tikzpicture}[scale=1.0, thick]
  \def\R{2.2}
  \draw (0,0) circle (\R);
  % 12 小时刻度
  \foreach \i in {1,...,12} {
    \pgfmathsetmacro{\ang}{90 - \i*30}
    \node at ({(\R-0.3)*cos(\ang)}, {(\R-0.3)*sin(\ang)}) {\footnotesize \i};
    \draw ({\R*cos(\ang)}, {\R*sin(\ang)}) -- ({(\R-0.1)*cos(\ang)}, {(\R-0.1)*sin(\ang)});
  }
  % 分针 极角 -30°
  \draw[-{Stealth[length=4pt]}, thick] (0,0) -- ({1.8*cos(-30)}, {1.8*sin(-30)});
  % 时针 极角 -10°（较短）
  \draw[-{Stealth[length=4pt]}, thick, red] (0,0) -- ({1.2*cos(-10)}, {1.2*sin(-10)});
  % 夹角弧
  \draw[red] (0.7,0) arc (-10:-30:0.7);
  \node[red, font=\small] at ({0.95*cos(-20)}, {0.95*sin(-20)}) {$20°$};
  \filldraw (0,0) circle (1.5pt);
  \node[below=2.5cm, font=\small] at (0,0) {3 时 20 分，时分针夹角 $20°$};
\end{tikzpicture}
\end{document}
